\daytitle{8}{功、功率与动能定理}

\begin{learningbox}{本讲义大致内容}
从力的角度解决问题需要知道加速度和运动时间，但很多问题我们只关心初末状态，不关心中间过程。能量方法就是为此而生。今天学习功、功率、动能定理——用能量"绕开"加速度，直接建立初末状态的关系。
\end{learningbox}

\section*{一、学习目标}

\begin{enumerate}
  \item 理解功的定义 $W = Fx\cos\theta$，会判断正功、负功、零功。
  \item 理解功率 $P = \dfrac{W}{t} = Fv$，区分瞬时功率和平均功率。
  \item 掌握动能 $E_k = \dfrac12 mv^2$。
  \item 掌握动能定理 $W_{\text{合}} = \Delta E_k$，会用动能定理解决多过程问题。
  \item 理解用能量方法比用动力学方法更简单的场景。
\end{enumerate}

\section*{二、初高中衔接}

初中已经学过"功 = 力 × 距离"，但有两个重要局限：
\begin{enumerate}
  \item 初中只考虑力与位移同向的情况，高中要考虑力的方向与位移方向有夹角的情况。
  \item 初中只算大小，不分正负。高中要区分正功（动力做功）和负功（阻力做功）。
\end{enumerate}

更重要的是：初中对"能量"只有定性概念。高中开始定量计算动能、势能，并用"功是能量转化的量度"这个核心思想来解题。

\section*{三、核心概念}

\begin{conceptbox}{功}

\textbf{定义：} 力对物体做的功等于力的大小、位移大小、力与位移夹角余弦的乘积。

\[
\boxed{W = F x \cos\theta}
\]

\begin{itemize}
  \item $F$：力的大小（$\mathrm{N}$）。
  \item $x$：位移的大小（$\mathrm{m}$）。
  \item $\theta$：力与位移方向的夹角。
\end{itemize}

\textbf{正功与负功：}
\begin{center}
\begin{tabular}{lll}
\toprule
$\theta$ 范围 & $\cos\theta$ & 功的性质 \\
\midrule
$0^\circ \leq \theta < 90^\circ$ & $>0$ & 正功（动力做功） \\
$\theta = 90^\circ$ & $=0$ & 不做功 \\
$90^\circ < \theta \leq 180^\circ$ & $<0$ & 负功（阻力做功） \\
\bottomrule
\end{tabular}
\end{center}

\begin{center}
% TikZ 图：力做功的示意图
% Day 08：功、功率与动能定理
\begin{tikzpicture}[scale=0.9,>=Latex]
  % 地面
  \draw[thick] (-1,0) -- (5,0);
  % 物体
  \draw[thick] (0,0) rectangle (1.5,1.2);
  \node at (0.75,0.6) {$m$};
  % 力 F（与水平成θ角）
  \draw[->,thick,red] (0.75,0.6) -- (2.8,2.0) node[above] {$F$};
  \draw[dashed,thick] (0.75,0.6) -- (2.5,0.6);
  \draw (1.2,0.6) arc (0:37:0.45);
  \node at (1.6,0.8) {$\theta$};
  % 位移
  \draw[->,thick,blue] (1.5,0.3) -- (4.5,0.3) node[below] {$x$};
  % 做功标注
  \node[below] at (2.5,-0.5) {$W = Fx\cos\theta$};
\end{tikzpicture}


\figcaption{力 $F$ 与位移 $x$ 夹角为 $\theta$ 时，$W = Fx\cos\theta$}\label{fig:work-done}
\end{center}

\textbf{常见注意：}
\begin{enumerate}
  \item 摩擦力可以做正功（传送带对物块）、负功（刹车）、不做功（静摩擦力的作用点无地面位移时，如人走路时鞋底静摩擦力对人不做功——但人整体动能增加是通过内部肌肉做功实现的）。
  \item 功是标量，没有方向，但有正负——正负表示能量是增加还是减少。
  \item 总功等于各力做功的代数和，不是矢量和。
\end{enumerate}

\end{conceptbox}

\begin{conceptbox}{功率}

\textbf{定义：} 功与完成这些功所用时间的比值。

\[
P = \frac{W}{t} \quad\text{(平均功率)}
\]
\[
P = F v \cos\theta \quad\text{(瞬时功率)}
\]

\begin{itemize}
  \item 单位：瓦特（$\mathrm{W}$），$1\ \mathrm{W} = 1\ \mathrm{J/s}$。
  \item 瞬时功率公式中 $v$ 是瞬时速度。
  \item 机车功率问题：$P = F v$，其中 $F$ 是牵引力。当功率一定时，牵引力与速度成反比。
\end{itemize}

\end{conceptbox}

\begin{conceptbox}{动能与动能定理}

\textbf{动能：}
\[
E_k = \frac12 m v^2
\]
\begin{itemize}
  \item 标量，单位焦耳（$\mathrm{J}$）。
  \item 动能非负（$v=0$ 时 $E_k=0$），与速度方向无关。
  \item 动能的增量只与初末速度有关，与中间过程无关。
\end{itemize}

\textbf{动能定理：}
\[
\boxed{W_{\text{合}} = \Delta E_k = \frac12 mv_2^2 - \frac12 mv_1^2}
\]

\textbf{来源（匀变速直线的一步推导）：}
若合力与位移同向且为恒力，
\[
W_{\text{合}} = F_{\text{合}}x = max
\]
又由 $v_2^2 - v_1^2 = 2ax$，得
\[
W_{\text{合}} = \frac12 mv_2^2 - \frac12 mv_1^2
\]
更一般的变力、曲线运动，可以把过程分成很多小段，每小段近似用上式，求和后仍得到
$W_{\text{合}}=\Delta E_k$。

\begin{itemize}
  \item $W_{\text{合}}$ 是合力做的功（或所有力做功的代数和）。
  \item 动能定理不涉及加速度和时间，适合处理多过程问题。
  \item 与牛顿定律 $F = ma$ 相比，动能定理是"从能量角度看运动"。
\end{itemize}

\begin{center}
% TikZ 图：动能定理示意
% Day 08：功、功率与动能定理
\begin{tikzpicture}[scale=0.9,>=Latex]
  \draw[thick] (-1,0) -- (7,0);
  % 起点
  \draw[thick] (0,0) rectangle (1.5,1.2);
  \node at (0.75,0.6) {$m$};
  \draw[->,thick,blue] (0.20,1.45) -- (1.30,1.45) node[midway,above] {$v_1$};
  % 终点
  \draw[thick] (5,0) rectangle (6.5,1.2);
  \node at (5.75,0.6) {$m$};
  \draw[->,thick,blue] (5.20,1.45) -- (6.30,1.45) node[midway,above] {$v_2$};
  % 恒定合力画在物块上，避免误看成连接两个状态的连线。
  \draw[->,ultra thick,red] (0.75,0.6) -- (2.25,0.6)
    node[midway,above] {$F_{\text{合}}$};
  \draw[->,thick,gray] (2.55,0.6) -- (4.45,0.6)
    node[midway,above,font=\small] {运动过程};
  % 位移
  \draw[<->] (0,-0.6) -- (6.5,-0.6) node[midway,below] {$x$};
  % 动能定理
  \node at (3.25,2.20) {$F_{\text{合}} x = \frac12 m v_2^2 - \frac12 m v_1^2$};
\end{tikzpicture}


\figcaption{动能定理：合外力做功等于动能变化}\label{fig:work-energy-theorem}
\end{center}

\end{conceptbox}

\section*{四、公式与适用条件}

\begin{center}
\begin{tabular}{lll}
\toprule
公式 & 物理意义 & 适用条件 \\
\midrule
$W = Fx\cos\theta$ & 恒力做功 & 恒力，位移已知 \\
$P = \dfrac{W}{t}$ & 平均功率 & 一段时间内的平均功率 \\
$P = Fv$ & 瞬时功率 & 力与速度同向时 \\
$E_k = \frac12 mv^2$ & 动能定义 & 任何运动 \\
$W_{\text{合}} = \Delta E_k$ & 动能定理 & 任何过程（变力也成立） \\
\bottomrule
\end{tabular}
\end{center}

\section*{五、典型模型}

\subsection*{模型 1：恒力做功}

直接用 $W = Fx\cos\theta$ 计算。重力做功 $W_G = mgh$ 只与高度差有关，与路径无关。

\subsection*{模型 2：变力做功}

若力随位移线性变化（如弹簧 $F = kx$），可用平均力 $\bar{F} = \frac{F_1+F_2}{2}$ 求功。

\textbf{弹簧弹力做功（区分两种视角）：}
\begin{itemize}
  \item \textbf{外力拉弹簧做正功：} $W_{\text{外}} = \frac12 kx^2$，弹簧弹性势能增加。
  \item \textbf{弹簧弹力做负功：} $W_{\text{弹}} = -\frac12 kx^2$，弹性势能增加。
\end{itemize}
两者大小相等、符号相反，但不能把它们相加成一个新的功。对弹簧来说，
\[
\Delta E_p = \frac12 kx^2,\qquad
W_{\text{外}}=\Delta E_p,\qquad
W_{\text{弹}}=-\Delta E_p
\]
下节课会详细展开弹性势能。

\subsection*{模型 3：机车启动}

恒定功率启动：$P = Fv$，随着 $v$ 增大、$F$ 减小，加速度减小。当 $F = f$（阻力）时，达到最大速度 $v_{\max} = \dfrac{P}{f}$。

\section*{六、例题精讲}

\begin{examplebox}{恒力做功的计算}

\textbf{题意分析：} 质量为 $m = 10\ \mathrm{kg}$ 的物体在水平面上，受与水平方向成 $37^\circ$ 角斜向上的拉力 $F = 50\ \mathrm{N}$ 作用，在水平面上前进了 $x = 10\ \mathrm{m}$，$\mu = 0.2$，$g = 10\ \mathrm{m/s^2}$。求各力做的功和总功。

\textbf{物理模型：} 恒力做功，受力分析后逐一计算。

\textbf{解题过程：}

拉力做功：$W_F = Fx\cos 37^\circ = 50 \times 10 \times 0.8 = 400\ \mathrm{J}$（正功）。

重力做功：$W_G = 0$（重力与位移方向垂直）。

支持力做功：$W_N = 0$（支持力与位移方向垂直）。

摩擦力：$N = mg - F\sin 37^\circ = 100 - 50 \times 0.6 = 70\ \mathrm{N}$，$f = \mu N = 0.2 \times 70 = 14\ \mathrm{N}$。

$W_f = f x \cos 180^\circ = 14 \times 10 \times (-1) = -140\ \mathrm{J}$（负功）。

总功：$W_{\text{合}} = 400 + 0 + 0 - 140 = 260\ \mathrm{J}$。

\textbf{方法总结：}
\begin{enumerate}
  \item 各个力分别算功，再代数相加得总功，不容易出错。
  \item 注意正压力 $N$ 不等于 $mg$——斜向上拉力会减小正压力。
\end{enumerate}

\end{examplebox}

\begin{examplebox}{动能定理的应用}

\textbf{题意分析：} 质量为 $m = 2\ \mathrm{kg}$ 的物体以 $v_0 = 6\ \mathrm{m/s}$ 的初速度沿粗糙水平面滑行，$\mu = 0.3$，$g = 10\ \mathrm{m/s^2}$。求物体能滑行的最大距离。

\textbf{物理模型：} 匀减速直到停止。用动能定理可以一步到位。

\textbf{解题过程：}

末速度 $v = 0$，合外力只有摩擦力做功：
\[
W_f = -\mu mg x
\]
由动能定理：
\[
-\mu mg x = 0 - \frac12 m v_0^2
\]
\[
x = \frac{v_0^2}{2\mu g} = \frac{6^2}{2 \times 0.3 \times 10} = \frac{36}{6} = 6\ \mathrm{m}
\]

\textbf{对比（用运动学）：} $a = -\mu g = -3\ \mathrm{m/s^2}$，$v^2 - v_0^2 = 2ax$，$0 - 36 = 2 \times (-3) \times x$，$x = 6\ \mathrm{m}$。结果相同，但动能定理避免了求加速度的中间步骤。

\textbf{方法总结：}
\begin{enumerate}
  \item 动能定理不涉及时间 $t$ 和加速度 $a$，直接从初末速度得到位移。
  \item 摩擦力做功 $W_f = -fx$，负号表示能量减少。
\end{enumerate}

\end{examplebox}

\begin{examplebox}{多过程动能定理}

\textbf{题意分析：} 物体从 $h = 5\ \mathrm{m}$ 高的光滑斜面顶端由静止滑下，然后在水平粗糙面上滑行，$\mu = 0.4$，$g = 10\ \mathrm{m/s^2}$。求物体在水平面上滑行的距离。

\textbf{物理模型：} 斜面（只有重力做功）+ 水平面（只有摩擦力做功）。用动能定理列全过程方程。

\textbf{解题过程：}

取全程（从斜面顶端到水平面停止）为研究过程：
\begin{itemize}
  \item 重力做功：$W_G = mgh$（沿斜面下滑，重力做正功）。
  \item 摩擦力做功：斜面光滑无摩擦，水平面 $W_f = -\mu mg x$（$x$ 为水平滑行距离）。
  \item 初 $v_1 = 0$，末 $v_2 = 0$，$\Delta E_k = 0$。
\end{itemize}

由动能定理：
\[
mgh - \mu mg x = 0
\]
\[
x = \frac{h}{\mu} = \frac{5}{0.4} = 12.5\ \mathrm{m}
\]

\textbf{方法总结：}
\begin{enumerate}
  \item 多过程问题用动能定理可以一次性列方程，不需要分段求解。
  \item 重力做功只与高度差有关，与具体路径无关——这是用能量方法的核心优势。
\end{enumerate}

\end{examplebox}

\section*{七、变式训练}

\begin{enumerate}
  \item 质量为 $m = 5\ \mathrm{kg}$ 的物体在水平拉力 $F = 20\ \mathrm{N}$ 作用下前进 $x = 4\ \mathrm{m}$，$\mu = 0.3$，$g = 10\ \mathrm{m/s^2}$。求各力做的功和总功。
  
  \item 一颗子弹以 $v_0 = 400\ \mathrm{m/s}$ 的速度射入固定的木块，射入深度 $d = 20\ \mathrm{cm}$。设阻力恒定，求子弹受到的平均阻力。若子弹质量 $m = 10\ \mathrm{g}$。
  
  \item 物体从 $h = 3.2\ \mathrm{m}$ 高处由静止沿光滑斜面滑下，然后在水平面上滑行 $x = 8\ \mathrm{m}$ 后停止。求水平面的动摩擦因数 $\mu$。
\end{enumerate}

\section*{八、易错点}

\begin{mistakebox}
\begin{enumerate}
  \item \textbf{误以为 $W = Fx$ 在所有情况下成立：} 必须有 $F$ 与 $x$ 同向。否则要乘 $\cos\theta$。垂直的情况做功为零。
  
  \item \textbf{混淆功率和机械效率：} 功率是做功快慢，效率是有用功与总功之比，两个概念不同。
  
  \item \textbf{动能定理中漏掉某个力的功：} 把所有力做的功加起来——重力、弹力、摩擦力、拉力，一个不能少。
  
  \item \textbf{用动能定理时初末动能搞反：} 末动能减初动能，不是初减末。
  
  \item \textbf{把变力做功误用 $Fx$ 公式：} $W = Fx$ 只适用于恒力。变力做功要用平均力法、图像面积法或功能关系。
\end{enumerate}
\end{mistakebox}

\section*{九、竞赛入门补充}

\begin{competitionbox}
\textbf{变力做功的微元求和思路}

当力随位置变化时，可以把位移分成很多小段 $\Delta x$，每段内力近似恒定，各段功相加取极限就是积分：
\[
W = \int_{x_1}^{x_2} F(x) \, dx
\]

在 $F$-$x$ 图上，变力做功等于图线与 $x$ 轴围成的面积。这是因为每一小段的微元功
$\Delta W \approx F(x)\Delta x$，对应一条窄矩形的面积；把所有窄矩形相加并取极限，
就得到曲线下方的面积。

\begin{center}
% TikZ 图：弹簧外力做功的 F-x 面积
% Day 08：功和动能定理
\begin{tikzpicture}[scale=0.9,>=Latex]
  \draw[->,thick] (0,0) -- (5,0) node[right] {$x$};
  \draw[->,thick] (0,0) -- (0,3.2) node[above] {$F_{\text{外}}$};

  \coordinate (O) at (0,0);
  \coordinate (A) at (4,0);
  \coordinate (B) at (4,2.4);

  \fill[gray!18] (O) -- (A) -- (B) -- cycle;
  \draw[thick] (O) -- (B) node[pos=0.65,above left] {$F=kx$};
  \draw[dashed] (B) -- (A);
  \draw[dashed] (B) -- (0,2.4);

  \node[below] at (A) {$x$};
  \node[left] at (0,2.4) {$kx$};
  \node at (2.55,0.85) {$W_{\text{外}}=\frac12 kx^2$};
  \node[below left] at (O) {$O$};
\end{tikzpicture}


\figcaption{线性变力做功：$F$-$x$ 图像下的面积就是功}\label{fig:spring-work-area}
\end{center}

\textbf{例题：} 弹簧的劲度系数 $k = 200\ \mathrm{N/m}$，从原长 $l_0$ 拉伸 $x = 0.1\ \mathrm{m}$，求外力克服弹力做的功。

\textbf{解：} 弹簧弹力 $F_{\text{弹}} = kx$ 随 $x$ 线性增大，方向指向原长（收缩方向）。外力 $F_{\text{外}}$ 与之等大反向，$F_{\text{外}} = kx$。

$F_{\text{外}}$-$x$ 图上，三角形面积即外力做的正功：
\[
W_{\text{外}} = \frac12 \times x \times kx = \frac12 k x^2 = \frac12 \times 200 \times 0.01 = 1\ \mathrm{J}
\]

弹簧弹力做负功：$W_{\text{弹}} = -W_{\text{外}} = -\frac12 kx^2 = -1\ \mathrm{J}$。

这两个结果分别对应弹性势能的变化：$W_{\text{外}} = \Delta E_p$，$W_{\text{弹}} = -\Delta E_p$。下节课会系统讲解弹性势能。
\end{competitionbox}

\section*{十、今日小结}

\begin{enumerate}
  \item 功 $W = Fx\cos\theta$，正功 $(\theta < 90^\circ)$、负功 $(\theta > 90^\circ)$、零功 $(\theta = 90^\circ)$。
  \item 动能 $E_k = \frac12 mv^2$，动能定理 $W_{\text{合}} = \Delta E_k$。
  \item 动能定理不涉及加速度和时间，解决多过程问题很高效。
  \item 重力做功只与高度差有关，与路径无关。
  \item 变力做功可用 $F$-$x$ 图面积或功能关系求解。
\end{enumerate}
